\begin{table*}[tb]
\centering
\small
\setlength{\tabcolsep}{6pt}
\begin{tabular}{l l c c c c}
\toprule
Axis & Judge & agrees with factor & called indistinguishable & pairs & inter-rater $\alpha$ \\
\midrule
Static stability (fBD) & human & -- (3 decided) & -- & 8 & 0.18 \\
 & vision--language & -- (1 decided) & 1/1 & 7 & -- \\
Motion persistence (MCFF-L) & human & 80\% {\footnotesize [40, 100]} & -- & 8 & 0.24 \\
 & vision--language & -- (2 decided) & -- & 3 & -- \\
\bottomrule
\end{tabular}
\caption{\textbf{Preliminary perceptual validation of the factor axes.} Each judge answers one question per axis and never which output is better. \emph{agree} is agreement with the factor's ordering on pairs the factor separates clearly, over pairs the judge decided; \emph{can't tell} is the share of near-tie pairs called indistinguishable, which is the correct answer there. The unit is the pair, majority-voted across raters, and the interval is a percentile bootstrap resampling pairs rather than responses, which are clustered within rater and pair. A dash marks a cell with too few decided pairs to carry a rate; the count is given instead. This is a preliminary check on a measurement instrument, not a powered study, and no system is ranked by it.}
\label{tab:perception}
\end{table*}